% =============================================================================
% 1.4 Intelligence artificielle et traitement du langage naturel
% =============================================================================
\section{Intelligence artificielle et traitement automatique du langage}
\label{sec:ch1-tal}

\subsection{Fondements du traitement automatique du langage}

Le traitement automatique du langage naturel est une discipline de
l'intelligence artificielle qui vise à faire analyser, comprendre et produire du
langage humain par des machines. Ses origines remontent aux années 1950, mais
c'est au cours de la dernière décennie qu'elle a connu des avancées décisives
grâce à l'apprentissage profond.

\subsubsection{Approches traditionnelles}

Les premières approches reposaient sur des règles linguistiques et des méthodes
statistiques~:

\begin{itemize}
  \item \btpterm{Analyse morphosyntaxique}~: segmentation en phrases,
        tokenisation, étiquetage morphologique~;
  \item \btpterm{Grammaires formelles}~: analyse syntaxique fondée sur des
        règles de production~;
  \item \btpterm{Modèles de Markov cachés}~: modélisation probabiliste de
        séquences d'états latents~;
  \item \btpterm{Pondération TF-IDF}~: mesure de l'importance d'un terme dans un
        document relativement à une collection.
\end{itemize}

Ces approches, efficaces sur des domaines restreints et des tâches bien
délimitées, peinaient à capturer la complexité et la polysémie du langage
naturel. Elles conservent néanmoins une pertinence pratique~: la recherche
lexicale, dont nous verrons l'importance en
section~\ref{sec:ch1-recherche}, en descend directement.

\subsubsection{L'ère de l'apprentissage profond}

L'avènement des réseaux de neurones profonds, à partir de 2012, a transformé le
domaine~:

\begin{itemize}
  \item \btpterm{Word2Vec} (Mikolov et coll., 2013)~: représentation
        vectorielle des mots capturant des régularités sémantiques par
        arithmétique vectorielle~;
  \item \btpterm{GloVe} (Pennington et coll., 2014)~: vecteurs de mots dérivés
        des statistiques globales de co-occurrence~;
  \item \btpterm{Réseaux récurrents et à mémoire longue} (Hochreiter \&
        Schmidhuber, 1997)~: traitement de séquences avec état interne~;
  \item \btpterm{Mécanisme d'attention et architecture Transformer} (Vaswani et
        coll., 2017)~: abandon de la récurrence au profit d'une mise en relation
        directe de toutes les positions d'une séquence.
\end{itemize}

\subsection{Les modèles de langage}

\subsubsection{Principe}

Un modèle de langage est un système capable d'estimer la probabilité d'une
séquence de mots et, par échantillonnage successif, de produire du texte. Les
modèles modernes, fondés sur l'architecture Transformer, sont entraînés sur des
corpus massifs et développent des représentations fines des régularités du
langage.

\begin{btpdef}[Grand modèle de langage]
Un grand modèle de langage est un modèle de langage neuronal dont le nombre de
paramètres et la taille du corpus d'entraînement sont suffisamment élevés pour
qu'apparaissent des capacités non explicitement entraînées, notamment le suivi
d'instructions et l'apprentissage à partir de quelques exemples fournis dans
l'invite.
\end{btpdef}

\subsubsection{Évolution chronologique}

\begin{btptable}{Jalons de l'évolution des modèles de langage}
\label{tab:jalons-llm}
\begin{tabularx}{\linewidth}{@{}C{1.3cm} L{2.3cm} Y L{2.5cm}@{}}
\toprule
\rowcolor{btppale}
\btpth{Année} & \btpth{Modèle} & \btpth{Apport principal} & \btpth{Ordre de grandeur} \\
\midrule
2018 & BERT & Contextualisation bidirectionnelle par masquage. & $3{,}4 \times 10^{8}$ \\
2019 & GPT-2 & Génération de texte long cohérent. & $1{,}5 \times 10^{9}$ \\
2020 & GPT-3 & Apprentissage à partir de quelques exemples. & $1{,}75 \times 10^{11}$ \\
2022 & Modèles alignés par instruction & Suivi fiable de consignes. & Comparable \\
2023 & Modèles multimodaux & Raisonnement étendu, entrée non textuelle. & Non publié \\
2023 & Modèles ouverts (famille Llama) & Diffusion des poids, exécution locale. & $7$ à $70 \times 10^{9}$ \\
2024 & Mélanges d'experts & Capacité accrue à coût d'inférence contenu. & Activation partielle \\
\bottomrule
\end{tabularx}
\end{btptable}

\subsubsection{Modèles multilingues}

Les modèles multilingues sont déterminants dans un contexte francophone comme
celui des appels d'offres camerounais. Les travaux de référence incluent
\btpterm{XLM-RoBERTa} (Conneau et coll., 2019), entraîné sur une centaine de
langues, \btpterm{mT5} (Xue et coll., 2020), variante multilingue d'un modèle
encodeur-décodeur, et les modèles de plongement multilingues à représentation
mixte dense et creuse, dont \btpterm{BGE-M3} (2023) est représentatif.

\subsubsection{Modèles spécialisés}

Le domaine des appels d'offres nécessite une compréhension précise de textes
juridiques, techniques et financiers. Des modèles adaptés à ces domaines ont
émergé, dont \btpterm{LegalBERT} pour les textes juridiques et
\btpterm{FinBERT} pour les textes financiers. Leur intérêt réside moins dans un
gain brut de performance que dans une meilleure calibration du vocabulaire
spécialisé, ce qui réduit les confusions terminologiques.

\begin{btpwarn}[Limite pour le contexte francophone africain]
Ces modèles spécialisés sont majoritairement entraînés sur des corpus
anglo-saxons, dont le droit et les pratiques contractuelles diffèrent
sensiblement du droit camerounais des marchés publics. Leur transposition n'est
donc pas immédiate, ce qui renforce l'intérêt d'une approche par récupération
dans un corpus documentaire local plutôt que par mémorisation paramétrique.
\end{btpwarn}

\subsection{Les représentations vectorielles}

\subsubsection{Définition et principes}

Une représentation vectorielle, ou plongement, est la projection d'un texte —
mot, phrase ou document — dans un espace vectoriel de dimension fixe. La
propriété recherchée est que la proximité géométrique traduise la proximité
sémantique~: deux textes de sens voisin sont représentés par des vecteurs
proches, ce qui permet de mesurer une similarité par simple calcul de distance.

\subsubsection{Évolution des approches}

La figure~\ref{fig:evolution-representations} retrace les cinq générations de
représentations textuelles.

\begin{figure}[htbp]
  \centering
  \resizebox{\linewidth}{!}{% Évolution des représentations vectorielles de texte
\begin{tikzpicture}[node distance=4mm]

  \node[btpboxplain, text width=11.2cm, align=left] (r1)
    {\textbf{1. Représentations creuses} \hfill \btplabel{TF-IDF, BM25}\\[1pt]
     \btplabel{\itshape Termes exacts, aucune compréhension sémantique}};
  \node[btpboxplain, text width=11.2cm, align=left, below=of r1] (r2)
    {\textbf{2. Plongements statiques} \hfill \btplabel{Word2Vec, GloVe}\\[1pt]
     \btplabel{\itshape Sémantique lexicale, mais insensible au contexte}};
  \node[btpbox, text width=11.2cm, align=left, below=of r2] (r3)
    {\textbf{3. Représentations contextuelles} \hfill \btplabel{ELMo, BERT}\\[1pt]
     \btplabel{\itshape Contexte bidirectionnel, gestion de la polysémie}};
  \node[btpboxalt, text width=11.2cm, align=left, below=of r3] (r4)
    {\textbf{4. Modèles mixtes dense + creux} \hfill \btplabel{SPLADE, BGE-M3}\\[1pt]
     \btplabel{\itshape Recherche hybride nativement supportée}};
  \node[btpboxalt, text width=11.2cm, align=left, below=of r4] (r5)
    {\textbf{5. Représentations multimodales} \hfill \btplabel{CLIP, ImageBind}\\[1pt]
     \btplabel{\itshape Texte, images, plans et schémas dans un espace commun}};

  \draw[btparrow, line width=1.6pt]
    ($(r1.west)+(-0.45,0.15)$) -- ($(r5.west)+(-0.45,-0.15)$)
    node[midway, rotate=90, anchor=south, btplabel, yshift=1pt]
    {Enrichissement sémantique};

\end{tikzpicture}
}
  \caption{Évolution des représentations vectorielles de texte}
  \label{fig:evolution-representations}
\end{figure}

\subsubsection{Représentations denses et creuses}

\begin{btptable}{Comparaison des représentations denses et creuses}
\label{tab:dense-creux}
\begin{tabularx}{\linewidth}{@{}L{1.7cm} L{3.4cm} Y Y@{}}
\toprule
\rowcolor{btppale}
\btpth{Type} & \btpth{Caractéristiques} & \btpth{Avantages} & \btpth{Inconvénients} \\
\midrule
Dense & Vecteur continu de dimension fixe (768 ou 1024 par exemple). & Capture la sémantique, tolère la reformulation, généralise bien. & Peine sur les termes rares, les identifiants et les références normatives. \\
Creux & Vecteur de la taille du vocabulaire, majoritairement nul. & Capture les termes exacts, directement interprétable. & Ne relie pas les synonymes ni les paraphrases. \\
\bottomrule
\end{tabularx}
\end{btptable}

Cette complémentarité n'est pas une nuance de second ordre~: elle fonde le choix
d'une recherche hybride, développé en section~\ref{sec:ch1-recherche}.

\subsection{Applications des grands modèles de langage}

Les capacités des grands modèles de langage recoupent notre problématique sur
plusieurs plans~:

\begin{enumerate}
  \item \btpterm{Génération de textes}~: rédaction de mémoires techniques, de
        notes méthodologiques et de sections administratives~;
  \item \btpterm{Extraction d'informations}~: identification des exigences,
        des entités nommées et des relations dans un dossier entrant~;
  \item \btpterm{Résumé}~: synthèse de pièces juridiques ou techniques longues~;
  \item \btpterm{Question-réponse documentaire}~: réponse à une interrogation
        précise à partir d'un corpus déterminé~;
  \item \btpterm{Classification}~: catégorisation automatique des pièces d'un
        dossier selon leur nature~;
  \item \btpterm{Reformulation de requêtes}~: enrichissement d'une question
        utilisateur avant recherche.
\end{enumerate}

Ces capacités partagent toutefois une faiblesse commune, dont la maîtrise
constitue le c\oe ur de notre travail~: en l'absence de sources fournies, un
modèle de langage produit un texte plausible, non un texte vrai. C'est cette
faiblesse que l'architecture examinée à la section~\ref{sec:ch1-rag} entend
corriger.
